% Source: https://artofproblemsolving.com/community/c3h1944535_mock_amc_8b_final_five
% Stable identity: aops-topic-1944535-q23
The tangent lengths from a vertex to the incircle are equal, and in a triangle with sides $u,v,w$ the tangent length from the vertex between sides $u$ and $v$ is $\dfrac{u+v-w}{2}$. Hence
\[
BP=\dfrac{BC+BX-CX}{2},\qquad BQ=\dfrac{BA+BX-AX}{2}.
\]
Since $AX=CX=40$,
\[
PQ=|BP-BQ|=\dfrac{|BC-BA|}{2}=\dfrac{64-48}{2}=\boxed{8},\qquad\textbf{(C)}.
\]

\begin{asy}
usepackage("amsmath");
size(5.5cm);
defaultpen(fontsize(10pt));
pair B=(0,0), C=(64,0), A=(0,48), X=(A+C)/2;
real bx=abs(X-B);
pair ic(pair P, pair Q, pair R){ real p=abs(Q-R), q=abs(R-P), r=abs(P-Q); return (p*P+q*Q+r*R)/(p+q+r); }
pair I1=ic(B,C,X), I2=ic(B,A,X);
real r1=abs(I1-(B+dot(I1-B,unit(X-B))*unit(X-B)));
real r2=abs(I2-(B+dot(I2-B,unit(X-B))*unit(X-B)));
pair P=B+dot(I1-B,unit(X-B))*unit(X-B), Q=B+dot(I2-B,unit(X-B))*unit(X-B);
draw(A--B--C--cycle,linewidth(1)); draw(B--X,linewidth(0.9));
draw(circle(I1,r1),blue); draw(circle(I2,r2),red);
draw((0,3)--(3,3)--(3,0));
dot(A);dot(B);dot(C);dot(X);dot(P,blue);dot(Q,red);
label("$A$",A,NW);label("$B$",B,SW);label("$C$",C,SE);label("$X$",X,NE);
label("$P$",P,SE,blue);label("$Q$",Q,NW,red);
label("$48$",(A+B)/2,W);label("$64$",(B+C)/2,S);
label("\mbox{}",(min(currentpicture,user=true)+max(currentpicture,user=true))/2,black); // keeps the text after the figure black when LaTeX typesets the labels inline
\end{asy}

(The triangle is right-angled at $B$ because $48^2+64^2=80^2$, so $BX=40$ as well, but this is not needed.)
